% !TEX root = ../mainfile.tex
% Event-study results tables. % !TEX root = ../mainfile.tex
% Event-study results tables. % !TEX root = ../mainfile.tex
% Event-study results tables. % !TEX root = ../mainfile.tex
% Event-study results tables. \input{content/tab_event_study_results} where needed,
% or copy the two table environments individually.
% Numbers generated by Code/03_event_study.py and Code/04_event_study_w50.py
% (constant-mean model, estimation window [-140,-21]; [-170,-51] for h=50).
% Verified against Data/processed/event_study/*.csv (re-run 2026-07-20).

\begin{table}[htbp]
\centering
\begin{threeparttable}
\caption{The U.S.\ dollar's response by event and horizon}
\label{tab:es_dollar_horizon}
\small
\setlength{\tabcolsep}{4.5pt}
\begin{tabular}{l rrr rrr}
\toprule
 & \multicolumn{6}{c}{Broad USD index --- CAR, \%} \\
\cmidrule(lr){2-7}
 & \multicolumn{3}{c}{Headline windows} & \multicolumn{3}{c}{Longer horizons (robustness)} \\
\cmidrule(lr){2-4}\cmidrule(lr){5-7}
Event & $(-1,+1)$ & $(0,+1)$ & $(0,+5)$ & $(0,+10)$ & $(0,+20)$ & $(0,+50)$ \\
\midrule
\multicolumn{7}{l}{\emph{Panel A. U.S.\ trade-policy shock}}\\
\quad Liberation Day        & $-1.70^{***}$ & $-1.44^{***}$ & $-0.10$       & $-3.00^{***}$ & $-3.98^{***}$ & $-6.34^{***}$ \\
\quad\quad Tariff pause     & $-1.48^{**}$  & $-1.48^{***}$ & $-3.07^{***}$ & $-3.44^{***}$ & $-4.51^{***}$ & $-6.74^{***}$ \\
\addlinespace
\multicolumn{7}{l}{\emph{Panel B. U.S.\ military / oil-supply shock}}\\
\quad Twelve-day war        & $-0.28$       & $+0.04$       & $+0.88$       & $-0.13$       & $+0.35$       & $-0.12$       \\
\quad Hormuz crisis         & $+1.34^{***}$ & $+1.39^{***}$ & $+1.53^{***}$ & $+2.10^{***}$ & $+3.25^{***}$ & $-0.13$       \\
\quad\quad Hormuz ceasefire & $-1.33^{***}$ & $-1.14^{***}$ & $-1.72^{***}$ & $-1.65^{**}$  & $-1.47$       & $-0.04$       \\
\addlinespace
\multicolumn{7}{l}{\emph{Panel C. External control (non-U.S.)}}\\
\quad Ukraine invasion      & $+0.45$       & $+0.49$       & $+1.03^{*}$   & $+1.65^{**}$  & $+0.84$       & $+3.12^{**}$  \\
\bottomrule
\end{tabular}
\begin{tablenotes}[flushleft]\footnotesize
\item \textit{Notes.} Cumulative abnormal returns of the broad nominal U.S.\ dollar
index (FRED \texttt{DTWEXBGS}), in percent, from a constant-mean model. Estimation
window $[-140,-21]$ trading days for $h\le 20$, ending where the event window begins;
$[-170,-51]$ for $h{=}50$ (pushed back so it does not overlap the wider event window).
Inference rests on the headline windows; the longer horizons gauge durability and
persistence but absorb confounding events and are read as robustness
(Section~\ref{sec:method_event}). Positive values denote dollar appreciation. Tariff
pause and Hormuz ceasefire are sub-events within their parent windows; ``Hormuz
ceasefire'' is the two-week ceasefire of 7~April 2026, distinct from the later 60-day
ceasefire of 11~June 2026. The US strikes (25~May 2026) and that 60-day June ceasefire
are omitted because the sample ends 30~June 2026, leaving too few post-event days for a
stable window.
$^{*}$, $^{**}$, $^{***}$ denote significance at the 10\%, 5\% and 1\% level.
\end{tablenotes}
\end{threeparttable}
\end{table}

\begin{table}[htbp]
\centering
\begin{threeparttable}
\caption{Cross-section of five-day responses: the mirror image}
\label{tab:es_cross_section}
\small
\setlength{\tabcolsep}{4.5pt}
\begin{tabular}{l rrr r}
\toprule
 & \multicolumn{3}{c}{U.S.-originated} & Control \\
\cmidrule(lr){2-4}\cmidrule(lr){5-5}
Series, $\CAR(0,5)$ (\%) & Liberation Day & Hormuz & Twelve-day war & Ukraine \\
\midrule
Broad USD index          & $-0.10$        & $+1.53^{***}$ & $+0.88$       & $+1.03^{*}$   \\
Dollar (equal-weighted)  & $+0.21$        & $+1.36^{**}$  & $+0.83$       & $+1.01$       \\
Safe currencies (carry)  & $+1.07$        & $-1.03$       & $-1.04$       & $-1.07^{*}$   \\
Risky currencies (carry) & $-0.53$        & $-1.97^{***}$ & $-1.03$       & $-2.56^{*}$   \\
Oil exporters            & $-0.80$        & $-0.69$       & $-0.45$       & $+0.30$       \\
Oil importers            & $+0.84$        & $-1.03^{*}$   & $-1.12$       & $-0.74$       \\
Gold                     & $-1.54$        & $-4.54$       & $-1.53$       & $+1.25$       \\
Euro Stoxx 50            & $-14.58^{***}$ & $-8.11^{***}$ & $-3.05$       & $-5.96^{**}$  \\
Brent crude              & $-14.32^{***}$ & $+27.63^{***}$& $+11.25^{***}$& $+13.82^{***}$\\
WTI crude                & $-13.14^{***}$ & $+34.55^{***}$& $+10.14^{**}$ & $+14.38^{***}$\\
\addlinespace
\multicolumn{5}{l}{\emph{Robustness: named safe-haven baskets}}\\
\quad Safe (JPY, CHF, EUR)                 & $+2.12^{*}$ & $-1.23$ & $-1.33$ & $-0.69$ \\
\quad Risky (TRY, BRL, ZAR, MXN, COP, IDR) & $-1.69^{*}$ & $-1.18$ & $-0.36$ & $+0.24$ \\
\addlinespace
\multicolumn{5}{l}{\emph{Robustness: euro-denominated commodities}}\\
\quad Brent crude in EUR & $-15.99^{***}$ & $+29.21^{***}$ & $+12.09^{**}$ & $+15.79^{***}$ \\
\quad Gold in EUR        & $-3.21$        & $-2.97$        & $-0.70$       & $+3.23^{*}$    \\
\bottomrule
\end{tabular}
\begin{tablenotes}[flushleft]\footnotesize
\item \textit{Notes.} Five-day cumulative abnormal returns (\%), constant-mean model,
estimation window $[-140,-21]$. Liberation Day is the trade-policy shock; Hormuz and the
twelve-day war are the military/oil-supply shocks; Ukraine is the non-U.S.\ control. For
the currency baskets a positive value means the basket \emph{appreciated against the
dollar}; for the dollar indices, gold and equities positive means a price increase, and
for crude positive means a price increase. \emph{Safe/Risky (carry)} are the bottom and
top terciles of the rebuilt forward-discount classification (Section~\ref{sec:data_portfolios});
the first robustness panel reports the literature-standard named-haven baskets. The
euro-denominated panel re-expresses the USD-priced commodities in euro
(Section~\ref{sec:data_fx}) to strip the mechanical dollar-in-the-denominator effect:
the crude response survives the change of numéraire, while gold's moves shrink toward
zero, consistent with gold trading as a dollar hedge. Gold's CARs
are baseline-sensitive given its exceptional 2025--26 trend and are read for sign rather
than magnitude (Section~\ref{sec:es_gold}). $^{*}$, $^{**}$, $^{***}$ denote significance
at the 10\%, 5\% and 1\% level.
\end{tablenotes}
\end{threeparttable}
\end{table}
 where needed,
% or copy the two table environments individually.
% Numbers generated by Code/03_event_study.py and Code/04_event_study_w50.py
% (constant-mean model, estimation window [-140,-21]; [-170,-51] for h=50).
% Verified against Data/processed/event_study/*.csv (re-run 2026-07-20).

\begin{table}[htbp]
\centering
\begin{threeparttable}
\caption{The U.S.\ dollar's response by event and horizon}
\label{tab:es_dollar_horizon}
\small
\setlength{\tabcolsep}{4.5pt}
\begin{tabular}{l rrr rrr}
\toprule
 & \multicolumn{6}{c}{Broad USD index --- CAR, \%} \\
\cmidrule(lr){2-7}
 & \multicolumn{3}{c}{Headline windows} & \multicolumn{3}{c}{Longer horizons (robustness)} \\
\cmidrule(lr){2-4}\cmidrule(lr){5-7}
Event & $(-1,+1)$ & $(0,+1)$ & $(0,+5)$ & $(0,+10)$ & $(0,+20)$ & $(0,+50)$ \\
\midrule
\multicolumn{7}{l}{\emph{Panel A. U.S.\ trade-policy shock}}\\
\quad Liberation Day        & $-1.70^{***}$ & $-1.44^{***}$ & $-0.10$       & $-3.00^{***}$ & $-3.98^{***}$ & $-6.34^{***}$ \\
\quad\quad Tariff pause     & $-1.48^{**}$  & $-1.48^{***}$ & $-3.07^{***}$ & $-3.44^{***}$ & $-4.51^{***}$ & $-6.74^{***}$ \\
\addlinespace
\multicolumn{7}{l}{\emph{Panel B. U.S.\ military / oil-supply shock}}\\
\quad Twelve-day war        & $-0.28$       & $+0.04$       & $+0.88$       & $-0.13$       & $+0.35$       & $-0.12$       \\
\quad Hormuz crisis         & $+1.34^{***}$ & $+1.39^{***}$ & $+1.53^{***}$ & $+2.10^{***}$ & $+3.25^{***}$ & $-0.13$       \\
\quad\quad Hormuz ceasefire & $-1.33^{***}$ & $-1.14^{***}$ & $-1.72^{***}$ & $-1.65^{**}$  & $-1.47$       & $-0.04$       \\
\addlinespace
\multicolumn{7}{l}{\emph{Panel C. External control (non-U.S.)}}\\
\quad Ukraine invasion      & $+0.45$       & $+0.49$       & $+1.03^{*}$   & $+1.65^{**}$  & $+0.84$       & $+3.12^{**}$  \\
\bottomrule
\end{tabular}
\begin{tablenotes}[flushleft]\footnotesize
\item \textit{Notes.} Cumulative abnormal returns of the broad nominal U.S.\ dollar
index (FRED \texttt{DTWEXBGS}), in percent, from a constant-mean model. Estimation
window $[-140,-21]$ trading days for $h\le 20$, ending where the event window begins;
$[-170,-51]$ for $h{=}50$ (pushed back so it does not overlap the wider event window).
Inference rests on the headline windows; the longer horizons gauge durability and
persistence but absorb confounding events and are read as robustness
(Section~\ref{sec:method_event}). Positive values denote dollar appreciation. Tariff
pause and Hormuz ceasefire are sub-events within their parent windows; ``Hormuz
ceasefire'' is the two-week ceasefire of 7~April 2026, distinct from the later 60-day
ceasefire of 11~June 2026. The US strikes (25~May 2026) and that 60-day June ceasefire
are omitted because the sample ends 30~June 2026, leaving too few post-event days for a
stable window.
$^{*}$, $^{**}$, $^{***}$ denote significance at the 10\%, 5\% and 1\% level.
\end{tablenotes}
\end{threeparttable}
\end{table}

\begin{table}[htbp]
\centering
\begin{threeparttable}
\caption{Cross-section of five-day responses: the mirror image}
\label{tab:es_cross_section}
\small
\setlength{\tabcolsep}{4.5pt}
\begin{tabular}{l rrr r}
\toprule
 & \multicolumn{3}{c}{U.S.-originated} & Control \\
\cmidrule(lr){2-4}\cmidrule(lr){5-5}
Series, $\CAR(0,5)$ (\%) & Liberation Day & Hormuz & Twelve-day war & Ukraine \\
\midrule
Broad USD index          & $-0.10$        & $+1.53^{***}$ & $+0.88$       & $+1.03^{*}$   \\
Dollar (equal-weighted)  & $+0.21$        & $+1.36^{**}$  & $+0.83$       & $+1.01$       \\
Safe currencies (carry)  & $+1.07$        & $-1.03$       & $-1.04$       & $-1.07^{*}$   \\
Risky currencies (carry) & $-0.53$        & $-1.97^{***}$ & $-1.03$       & $-2.56^{*}$   \\
Oil exporters            & $-0.80$        & $-0.69$       & $-0.45$       & $+0.30$       \\
Oil importers            & $+0.84$        & $-1.03^{*}$   & $-1.12$       & $-0.74$       \\
Gold                     & $-1.54$        & $-4.54$       & $-1.53$       & $+1.25$       \\
Euro Stoxx 50            & $-14.58^{***}$ & $-8.11^{***}$ & $-3.05$       & $-5.96^{**}$  \\
Brent crude              & $-14.32^{***}$ & $+27.63^{***}$& $+11.25^{***}$& $+13.82^{***}$\\
WTI crude                & $-13.14^{***}$ & $+34.55^{***}$& $+10.14^{**}$ & $+14.38^{***}$\\
\addlinespace
\multicolumn{5}{l}{\emph{Robustness: named safe-haven baskets}}\\
\quad Safe (JPY, CHF, EUR)                 & $+2.12^{*}$ & $-1.23$ & $-1.33$ & $-0.69$ \\
\quad Risky (TRY, BRL, ZAR, MXN, COP, IDR) & $-1.69^{*}$ & $-1.18$ & $-0.36$ & $+0.24$ \\
\addlinespace
\multicolumn{5}{l}{\emph{Robustness: euro-denominated commodities}}\\
\quad Brent crude in EUR & $-15.99^{***}$ & $+29.21^{***}$ & $+12.09^{**}$ & $+15.79^{***}$ \\
\quad Gold in EUR        & $-3.21$        & $-2.97$        & $-0.70$       & $+3.23^{*}$    \\
\bottomrule
\end{tabular}
\begin{tablenotes}[flushleft]\footnotesize
\item \textit{Notes.} Five-day cumulative abnormal returns (\%), constant-mean model,
estimation window $[-140,-21]$. Liberation Day is the trade-policy shock; Hormuz and the
twelve-day war are the military/oil-supply shocks; Ukraine is the non-U.S.\ control. For
the currency baskets a positive value means the basket \emph{appreciated against the
dollar}; for the dollar indices, gold and equities positive means a price increase, and
for crude positive means a price increase. \emph{Safe/Risky (carry)} are the bottom and
top terciles of the rebuilt forward-discount classification (Section~\ref{sec:data_portfolios});
the first robustness panel reports the literature-standard named-haven baskets. The
euro-denominated panel re-expresses the USD-priced commodities in euro
(Section~\ref{sec:data_fx}) to strip the mechanical dollar-in-the-denominator effect:
the crude response survives the change of numéraire, while gold's moves shrink toward
zero, consistent with gold trading as a dollar hedge. Gold's CARs
are baseline-sensitive given its exceptional 2025--26 trend and are read for sign rather
than magnitude (Section~\ref{sec:es_gold}). $^{*}$, $^{**}$, $^{***}$ denote significance
at the 10\%, 5\% and 1\% level.
\end{tablenotes}
\end{threeparttable}
\end{table}
 where needed,
% or copy the two table environments individually.
% Numbers generated by Code/03_event_study.py and Code/04_event_study_w50.py
% (constant-mean model, estimation window [-140,-21]; [-170,-51] for h=50).
% Verified against Data/processed/event_study/*.csv (re-run 2026-07-20).

\begin{table}[htbp]
\centering
\begin{threeparttable}
\caption{The U.S.\ dollar's response by event and horizon}
\label{tab:es_dollar_horizon}
\small
\setlength{\tabcolsep}{4.5pt}
\begin{tabular}{l rrr rrr}
\toprule
 & \multicolumn{6}{c}{Broad USD index --- CAR, \%} \\
\cmidrule(lr){2-7}
 & \multicolumn{3}{c}{Headline windows} & \multicolumn{3}{c}{Longer horizons (robustness)} \\
\cmidrule(lr){2-4}\cmidrule(lr){5-7}
Event & $(-1,+1)$ & $(0,+1)$ & $(0,+5)$ & $(0,+10)$ & $(0,+20)$ & $(0,+50)$ \\
\midrule
\multicolumn{7}{l}{\emph{Panel A. U.S.\ trade-policy shock}}\\
\quad Liberation Day        & $-1.70^{***}$ & $-1.44^{***}$ & $-0.10$       & $-3.00^{***}$ & $-3.98^{***}$ & $-6.34^{***}$ \\
\quad\quad Tariff pause     & $-1.48^{**}$  & $-1.48^{***}$ & $-3.07^{***}$ & $-3.44^{***}$ & $-4.51^{***}$ & $-6.74^{***}$ \\
\addlinespace
\multicolumn{7}{l}{\emph{Panel B. U.S.\ military / oil-supply shock}}\\
\quad Twelve-day war        & $-0.28$       & $+0.04$       & $+0.88$       & $-0.13$       & $+0.35$       & $-0.12$       \\
\quad Hormuz crisis         & $+1.34^{***}$ & $+1.39^{***}$ & $+1.53^{***}$ & $+2.10^{***}$ & $+3.25^{***}$ & $-0.13$       \\
\quad\quad Hormuz ceasefire & $-1.33^{***}$ & $-1.14^{***}$ & $-1.72^{***}$ & $-1.65^{**}$  & $-1.47$       & $-0.04$       \\
\addlinespace
\multicolumn{7}{l}{\emph{Panel C. External control (non-U.S.)}}\\
\quad Ukraine invasion      & $+0.45$       & $+0.49$       & $+1.03^{*}$   & $+1.65^{**}$  & $+0.84$       & $+3.12^{**}$  \\
\bottomrule
\end{tabular}
\begin{tablenotes}[flushleft]\footnotesize
\item \textit{Notes.} Cumulative abnormal returns of the broad nominal U.S.\ dollar
index (FRED \texttt{DTWEXBGS}), in percent, from a constant-mean model. Estimation
window $[-140,-21]$ trading days for $h\le 20$, ending where the event window begins;
$[-170,-51]$ for $h{=}50$ (pushed back so it does not overlap the wider event window).
Inference rests on the headline windows; the longer horizons gauge durability and
persistence but absorb confounding events and are read as robustness
(Section~\ref{sec:method_event}). Positive values denote dollar appreciation. Tariff
pause and Hormuz ceasefire are sub-events within their parent windows; ``Hormuz
ceasefire'' is the two-week ceasefire of 7~April 2026, distinct from the later 60-day
ceasefire of 11~June 2026. The US strikes (25~May 2026) and that 60-day June ceasefire
are omitted because the sample ends 30~June 2026, leaving too few post-event days for a
stable window.
$^{*}$, $^{**}$, $^{***}$ denote significance at the 10\%, 5\% and 1\% level.
\end{tablenotes}
\end{threeparttable}
\end{table}

\begin{table}[htbp]
\centering
\begin{threeparttable}
\caption{Cross-section of five-day responses: the mirror image}
\label{tab:es_cross_section}
\small
\setlength{\tabcolsep}{4.5pt}
\begin{tabular}{l rrr r}
\toprule
 & \multicolumn{3}{c}{U.S.-originated} & Control \\
\cmidrule(lr){2-4}\cmidrule(lr){5-5}
Series, $\CAR(0,5)$ (\%) & Liberation Day & Hormuz & Twelve-day war & Ukraine \\
\midrule
Broad USD index          & $-0.10$        & $+1.53^{***}$ & $+0.88$       & $+1.03^{*}$   \\
Dollar (equal-weighted)  & $+0.21$        & $+1.36^{**}$  & $+0.83$       & $+1.01$       \\
Safe currencies (carry)  & $+1.07$        & $-1.03$       & $-1.04$       & $-1.07^{*}$   \\
Risky currencies (carry) & $-0.53$        & $-1.97^{***}$ & $-1.03$       & $-2.56^{*}$   \\
Oil exporters            & $-0.80$        & $-0.69$       & $-0.45$       & $+0.30$       \\
Oil importers            & $+0.84$        & $-1.03^{*}$   & $-1.12$       & $-0.74$       \\
Gold                     & $-1.54$        & $-4.54$       & $-1.53$       & $+1.25$       \\
Euro Stoxx 50            & $-14.58^{***}$ & $-8.11^{***}$ & $-3.05$       & $-5.96^{**}$  \\
Brent crude              & $-14.32^{***}$ & $+27.63^{***}$& $+11.25^{***}$& $+13.82^{***}$\\
WTI crude                & $-13.14^{***}$ & $+34.55^{***}$& $+10.14^{**}$ & $+14.38^{***}$\\
\addlinespace
\multicolumn{5}{l}{\emph{Robustness: named safe-haven baskets}}\\
\quad Safe (JPY, CHF, EUR)                 & $+2.12^{*}$ & $-1.23$ & $-1.33$ & $-0.69$ \\
\quad Risky (TRY, BRL, ZAR, MXN, COP, IDR) & $-1.69^{*}$ & $-1.18$ & $-0.36$ & $+0.24$ \\
\addlinespace
\multicolumn{5}{l}{\emph{Robustness: euro-denominated commodities}}\\
\quad Brent crude in EUR & $-15.99^{***}$ & $+29.21^{***}$ & $+12.09^{**}$ & $+15.79^{***}$ \\
\quad Gold in EUR        & $-3.21$        & $-2.97$        & $-0.70$       & $+3.23^{*}$    \\
\bottomrule
\end{tabular}
\begin{tablenotes}[flushleft]\footnotesize
\item \textit{Notes.} Five-day cumulative abnormal returns (\%), constant-mean model,
estimation window $[-140,-21]$. Liberation Day is the trade-policy shock; Hormuz and the
twelve-day war are the military/oil-supply shocks; Ukraine is the non-U.S.\ control. For
the currency baskets a positive value means the basket \emph{appreciated against the
dollar}; for the dollar indices, gold and equities positive means a price increase, and
for crude positive means a price increase. \emph{Safe/Risky (carry)} are the bottom and
top terciles of the rebuilt forward-discount classification (Section~\ref{sec:data_portfolios});
the first robustness panel reports the literature-standard named-haven baskets. The
euro-denominated panel re-expresses the USD-priced commodities in euro
(Section~\ref{sec:data_fx}) to strip the mechanical dollar-in-the-denominator effect:
the crude response survives the change of numéraire, while gold's moves shrink toward
zero, consistent with gold trading as a dollar hedge. Gold's CARs
are baseline-sensitive given its exceptional 2025--26 trend and are read for sign rather
than magnitude (Section~\ref{sec:es_gold}). $^{*}$, $^{**}$, $^{***}$ denote significance
at the 10\%, 5\% and 1\% level.
\end{tablenotes}
\end{threeparttable}
\end{table}
 where needed,
% or copy the two table environments individually.
% Numbers generated by Code/03_event_study.py and Code/04_event_study_w50.py
% (constant-mean model, estimation window [-140,-21]; [-170,-51] for h=50).
% Verified against Data/processed/event_study/*.csv (re-run 2026-07-20).

\begin{table}[htbp]
\centering
\begin{threeparttable}
\caption{The U.S.\ dollar's response by event and horizon}
\label{tab:es_dollar_horizon}
\small
\setlength{\tabcolsep}{4.5pt}
\begin{tabular}{l rrr rrr}
\toprule
 & \multicolumn{6}{c}{Broad USD index --- CAR, \%} \\
\cmidrule(lr){2-7}
 & \multicolumn{3}{c}{Headline windows} & \multicolumn{3}{c}{Longer horizons (robustness)} \\
\cmidrule(lr){2-4}\cmidrule(lr){5-7}
Event & $(-1,+1)$ & $(0,+1)$ & $(0,+5)$ & $(0,+10)$ & $(0,+20)$ & $(0,+50)$ \\
\midrule
\multicolumn{7}{l}{\emph{Panel A. U.S.\ trade-policy shock}}\\
\quad Liberation Day        & $-1.70^{***}$ & $-1.44^{***}$ & $-0.10$       & $-3.00^{***}$ & $-3.98^{***}$ & $-6.34^{***}$ \\
\quad\quad Tariff pause     & $-1.48^{**}$  & $-1.48^{***}$ & $-3.07^{***}$ & $-3.44^{***}$ & $-4.51^{***}$ & $-6.74^{***}$ \\
\addlinespace
\multicolumn{7}{l}{\emph{Panel B. U.S.\ military / oil-supply shock}}\\
\quad Twelve-day war        & $-0.28$       & $+0.04$       & $+0.88$       & $-0.13$       & $+0.35$       & $-0.12$       \\
\quad Hormuz crisis         & $+1.34^{***}$ & $+1.39^{***}$ & $+1.53^{***}$ & $+2.10^{***}$ & $+3.25^{***}$ & $-0.13$       \\
\quad\quad Hormuz ceasefire & $-1.33^{***}$ & $-1.14^{***}$ & $-1.72^{***}$ & $-1.65^{**}$  & $-1.47$       & $-0.04$       \\
\addlinespace
\multicolumn{7}{l}{\emph{Panel C. External control (non-U.S.)}}\\
\quad Ukraine invasion      & $+0.45$       & $+0.49$       & $+1.03^{*}$   & $+1.65^{**}$  & $+0.84$       & $+3.12^{**}$  \\
\bottomrule
\end{tabular}
\begin{tablenotes}[flushleft]\footnotesize
\item \textit{Notes.} Cumulative abnormal returns of the broad nominal U.S.\ dollar
index (FRED \texttt{DTWEXBGS}), in percent, from a constant-mean model. Estimation
window $[-140,-21]$ trading days for $h\le 20$, ending where the event window begins;
$[-170,-51]$ for $h{=}50$ (pushed back so it does not overlap the wider event window).
Inference rests on the headline windows; the longer horizons gauge durability and
persistence but absorb confounding events and are read as robustness
(Section~\ref{sec:method_event}). Positive values denote dollar appreciation. Tariff
pause and Hormuz ceasefire are sub-events within their parent windows; ``Hormuz
ceasefire'' is the two-week ceasefire of 7~April 2026, distinct from the later 60-day
ceasefire of 11~June 2026. The US strikes (25~May 2026) and that 60-day June ceasefire
are omitted because the sample ends 30~June 2026, leaving too few post-event days for a
stable window.
$^{*}$, $^{**}$, $^{***}$ denote significance at the 10\%, 5\% and 1\% level.
\end{tablenotes}
\end{threeparttable}
\end{table}

\begin{table}[htbp]
\centering
\begin{threeparttable}
\caption{Cross-section of five-day responses: the mirror image}
\label{tab:es_cross_section}
\small
\setlength{\tabcolsep}{4.5pt}
\begin{tabular}{l rrr r}
\toprule
 & \multicolumn{3}{c}{U.S.-originated} & Control \\
\cmidrule(lr){2-4}\cmidrule(lr){5-5}
Series, $\CAR(0,5)$ (\%) & Liberation Day & Hormuz & Twelve-day war & Ukraine \\
\midrule
Broad USD index          & $-0.10$        & $+1.53^{***}$ & $+0.88$       & $+1.03^{*}$   \\
Dollar (equal-weighted)  & $+0.21$        & $+1.36^{**}$  & $+0.83$       & $+1.01$       \\
Safe currencies (carry)  & $+1.07$        & $-1.03$       & $-1.04$       & $-1.07^{*}$   \\
Risky currencies (carry) & $-0.53$        & $-1.97^{***}$ & $-1.03$       & $-2.56^{*}$   \\
Oil exporters            & $-0.80$        & $-0.69$       & $-0.45$       & $+0.30$       \\
Oil importers            & $+0.84$        & $-1.03^{*}$   & $-1.12$       & $-0.74$       \\
Gold                     & $-1.54$        & $-4.54$       & $-1.53$       & $+1.25$       \\
Euro Stoxx 50            & $-14.58^{***}$ & $-8.11^{***}$ & $-3.05$       & $-5.96^{**}$  \\
Brent crude              & $-14.32^{***}$ & $+27.63^{***}$& $+11.25^{***}$& $+13.82^{***}$\\
WTI crude                & $-13.14^{***}$ & $+34.55^{***}$& $+10.14^{**}$ & $+14.38^{***}$\\
\addlinespace
\multicolumn{5}{l}{\emph{Robustness: named safe-haven baskets}}\\
\quad Safe (JPY, CHF, EUR)                 & $+2.12^{*}$ & $-1.23$ & $-1.33$ & $-0.69$ \\
\quad Risky (TRY, BRL, ZAR, MXN, COP, IDR) & $-1.69^{*}$ & $-1.18$ & $-0.36$ & $+0.24$ \\
\addlinespace
\multicolumn{5}{l}{\emph{Robustness: euro-denominated commodities}}\\
\quad Brent crude in EUR & $-15.99^{***}$ & $+29.21^{***}$ & $+12.09^{**}$ & $+15.79^{***}$ \\
\quad Gold in EUR        & $-3.21$        & $-2.97$        & $-0.70$       & $+3.23^{*}$    \\
\bottomrule
\end{tabular}
\begin{tablenotes}[flushleft]\footnotesize
\item \textit{Notes.} Five-day cumulative abnormal returns (\%), constant-mean model,
estimation window $[-140,-21]$. Liberation Day is the trade-policy shock; Hormuz and the
twelve-day war are the military/oil-supply shocks; Ukraine is the non-U.S.\ control. For
the currency baskets a positive value means the basket \emph{appreciated against the
dollar}; for the dollar indices, gold and equities positive means a price increase, and
for crude positive means a price increase. \emph{Safe/Risky (carry)} are the bottom and
top terciles of the rebuilt forward-discount classification (Section~\ref{sec:data_portfolios});
the first robustness panel reports the literature-standard named-haven baskets. The
euro-denominated panel re-expresses the USD-priced commodities in euro
(Section~\ref{sec:data_fx}) to strip the mechanical dollar-in-the-denominator effect:
the crude response survives the change of numéraire, while gold's moves shrink toward
zero, consistent with gold trading as a dollar hedge. Gold's CARs
are baseline-sensitive given its exceptional 2025--26 trend and are read for sign rather
than magnitude (Section~\ref{sec:es_gold}). $^{*}$, $^{**}$, $^{***}$ denote significance
at the 10\%, 5\% and 1\% level.
\end{tablenotes}
\end{threeparttable}
\end{table}
